\subsection{DCache 集成（exp22）}

将 Cache 模块作为 DCache 集成到 CPU 访存通路中，需要处理 EX 级与 MEM 级的握手协议、非缓存访问的统一处理，以及强序非缓存的顺序性保证。

\subsubsection{DCache 实例化与接口连接}

\begin{lstlisting}[language=Verilog, caption=DCache 实例化]
cache #(
    .NUM_WAYS(2)
) u_dcache (
    .clk(clk),
    .resetn(resetn),
    
    // CPU 侧接口
    .valid(data_sram_req),
    .op(data_sram_wr),            // 1: 写, 0: 读
    .index(data_vaddr[11:4]),     // 虚地址 [11:4] 作为 Index
    .tag(data_paddr[31:12]),      // 物理地址 [31:12] 作为 Tag
    .offset(data_vaddr[3:0]),
    .wstrb(data_sram_wstrb),
    .wdata(data_sram_wdata),
    .cacheable(data_mat != 2'b10),  // MAT != 2'b10 表示可缓存
    .addr_ok(data_sram_addr_ok),
    .data_ok(data_sram_data_ok),
    .rdata(data_sram_rdata),
    .cacop(ex_inst_dcache_op),      // CACOP 操作 DCache
    .cacop_mode(ex_cacop_mode),
    
    // AXI 侧接口
    .rd_req(dcache_rd_req),
    .rd_type(dcache_rd_type),
    .rd_addr(dcache_rd_addr),
    .rd_rdy(dcache_rd_rdy),
    .ret_valid(dcache_ret_valid),
    .ret_last(dcache_ret_last),
    .ret_data(dcache_ret_data),
    .wr_req(dcache_wr_req),
    .wr_type(dcache_wr_type),
    .wr_addr(dcache_wr_addr),
    .wr_data(dcache_wr_data),
    .wr_wstrb(dcache_wr_wstrb),
    .wr_rdy(dcache_wr_rdy)
);
\end{lstlisting}

与 ICache 的主要区别：
\begin{itemize}
    \item \textbf{支持读写}：\verb|op| 信号区分读写操作，写操作需要 \verb|wstrb| 和 \verb|wdata|
    \item \textbf{需要写通道}：DCache 会发起写请求，\verb|wr_rdy| 来自 AXI Bridge
    \item \textbf{CACOP 来自 EX 级}：数据 Cache 操作在 EX 级进行
\end{itemize}

\subsubsection{访存地址翻译}

\begin{lstlisting}[language=Verilog, caption=访存地址翻译逻辑]
// TLB 查找（EX 级）
wire [18:0] ex_s1_vppn = data_vaddr[31:13];
wire        ex_s1_va_bit12 = data_vaddr[12];

tlb_searcher #(.TLBNUM(TLBNUM)) u_tlb_s1 (
    // TLB 表项输入（省略...）
    .vppn(ex_s1_vppn),
    .va_bit12(ex_s1_va_bit12),
    .asid(csr_asid[`CSR_ASID_ASID]),
    
    .found(ex_s1_found),
    .index(ex_s1_index),
    .ppn(ex_s1_ppn),
    .ps(ex_s1_ps),
    .plv(ex_s1_plv),
    .mat(ex_s1_mat),
    .d(ex_s1_d),
    .v(ex_s1_v)
);

// 三级地址翻译
assign ex_da = csr_crmd[`CSR_CRMD_DA] && !csr_crmd[`CSR_CRMD_PG];
assign ex_dmw0_hit = (data_vaddr[31:29] == csr_dmw0[31:29]) && 
                     dmw_plv_allowed(csr_dmw0, csr_crmd[`CSR_CRMD_PLV]);
assign ex_dmw1_hit = (data_vaddr[31:29] == csr_dmw1[31:29]) && 
                     dmw_plv_allowed(csr_dmw1, csr_crmd[`CSR_CRMD_PLV]);

// 物理地址与访问类型生成
assign data_paddr = ex_da ? data_vaddr :
                    ex_dmw0_hit ? {csr_dmw0[`CSR_DMW_PSEG], data_vaddr[27:0]} :
                    ex_dmw1_hit ? {csr_dmw1[`CSR_DMW_PSEG], data_vaddr[27:0]} :
                    ex_s1_found ? {ex_s1_ppn, data_vaddr[11:0]} :
                    data_vaddr;

assign data_mat = ex_da ? csr_crmd[`CSR_CRMD_DATM] :
                  ex_dmw0_hit ? csr_dmw0[`CSR_DMW_MAT] :
                  ex_dmw1_hit ? csr_dmw1[`CSR_DMW_MAT] :
                  ex_s1_found ? ex_s1_mat :
                  2'b10;

// 访存异常检测
wire ex_tlb_refill = !ex_da && !ex_dmw0_hit && !ex_dmw1_hit && !ex_s1_found;
wire ex_page_invalid = ex_s1_found && !ex_s1_v;
wire ex_page_dirty = is_store && ex_s1_found && !ex_s1_d;
wire ex_page_priv = ex_s1_found && (ex_s1_plv < csr_crmd[`CSR_CRMD_PLV]);
\end{lstlisting}

与取指地址翻译的区别：
\begin{itemize}
    \item \textbf{DA 模式使用 DATM}：直接地址翻译模式下，访存使用 DATM 域
    \item \textbf{检查 D 位}：写操作需要检查 Dirty 位，D=0 时触发 Page Dirty 异常
    \item \textbf{检查特权等级}：页表项的 PLV 不能低于当前 PLV
\end{itemize}

\subsubsection{流水线控制}

\begin{lstlisting}[language=Verilog, caption=EX 和 MEM 级控制逻辑]
// EX 级：发起访存请求
wire is_mem_op = is_store || is_load || ex_inst_dcache_op;

assign data_sram_req = ex_validout && is_mem_op && !ex_ex_valid;
assign data_sram_wr = is_store;
assign data_sram_addr = data_vaddr;
assign data_sram_wstrb = ex_mem_wstrb;
assign data_sram_wdata = ex_mem_wdata;

// MEM 级：接收访存数据
assign mem_allowin = !mem_validout || mem_allowout;
assign mem_allowout = wb_allowin || !mem_validout;

// 请求地址被接收后可以流水
assign ex_allowout = (mem_allowin && !is_mem_op) ||
                     (mem_allowin && is_mem_op && data_sram_addr_ok);

always @(posedge clk) begin
    if (rst || flush) begin
        mem_validout <= 1'b0;
    end else if (mem_allowin) begin
        mem_validout <= ex_validout && ex_allowout;
    end
end

// 数据返回后可以流水
assign mem_ready_go = !is_mem_op_r || data_sram_data_ok;
assign mem_allowout = mem_ready_go && wb_allowin;

// MEM 级数据寄存
always @(posedge clk) begin
    if (rst) begin
        mem_data <= 32'h0;
    end else if (data_sram_data_ok && mem_validout) begin
        mem_data <= data_sram_rdata;
    end
end
\end{lstlisting}

控制流关键点：
\begin{itemize}
    \item \textbf{地址握手与流水}：EX 级在 \verb|addr_ok=1| 时可流入 MEM 级
    \item \textbf{数据等待}：MEM 级在 \verb|data_ok=1| 时可流入 WB 级
    \item \textbf{写操作放行}：写操作的 \verb|data_ok| 在 LOOKUP 阶段即置 1，不阻塞流水
    \item \textbf{CACOP 处理}：CACOP 操作 DCache 视为访存操作，占用访存端口
\end{itemize}

\subsubsection{非缓存访问处理}

\begin{lstlisting}[language=Verilog, caption=Cache 模块中的非缓存访问处理]
// 在 Cache 模块内部
wire uncached_load = !buf_isstore && !buf_cacheable;
wire uncached_store = buf_isstore && !buf_cacheable;

// 非缓存访问强制视为 Cache Miss
wire effective_hit = hit && buf_cacheable && !(cacop_mode0 || ...);

// 状态机：非缓存写在 REPLACE 状态完成
MAIN_REPLACE: begin
    if (uncached_store) begin
        if (!wr_rdy)
            main_next_state = MAIN_REPLACE;  // 等待写就绪
        else
            main_next_state = MAIN_IDLE;     // 写完成，直接返回
    end else begin
        if (!rd_rdy)
            main_next_state = MAIN_REPLACE;
        else
            main_next_state = MAIN_REFILL;
    end
end

// 非缓存写的 data_ok
wire uncached_store_ok = (main_state == MAIN_REPLACE) && 
                         buf_isstore && !buf_cacheable && wr_rdy;

// 非缓存读在 REFILL 第一拍返回
wire refill_data_ok = (main_state == MAIN_REFILL) && ret_valid && 
    ~buf_isstore && (buf_cacheable ? (rpbuf_numrecv == buf_bank) : 1'b1);
\end{lstlisting}

非缓存访问处理要点：
\begin{itemize}
    \item \textbf{强制 Miss}：\verb|effective_hit=0|，进入 Miss 处理流程
    \item \textbf{不替换 Cache}：虽然走 Miss 流程，但不真正读写 Cache RAM
    \item \textbf{写完即返回}：非缓存写在 REPLACE 状态发出写请求后直接返回 IDLE
    \item \textbf{读第一拍返回}：非缓存读只需一个字，第一拍数据返回即完成
\end{itemize}

\subsubsection{强序非缓存的顺序性保证}

AXI Bridge 需要确保非缓存访问的严格顺序：

\begin{lstlisting}[language=Verilog, caption=AXI Bridge 非缓存顺序性控制]
// 非缓存写标记
reg uncached_write_pending;

always @(posedge clk) begin
    if (!resetn) begin
        uncached_write_pending <= 1'b0;
    end else begin
        if (dcache_wr_req && dcache_wr_type != 3'b100 && wr_rdy)
            uncached_write_pending <= 1'b1;  // 非缓存写开始
        if (bvalid && bready)
            uncached_write_pending <= 1'b0;  // bvalid 返回，写完成
    end
end

// 阻塞后续非缓存请求
assign dcache_rd_rdy = (!uncached_write_pending || 
                        dcache_rd_type == 3'b100) && r_state == R_IDLE;
assign dcache_wr_rdy = (!uncached_write_pending || 
                        dcache_wr_type == 3'b100) && wr_buf_empty;
\end{lstlisting}

顺序性保证机制：
\begin{itemize}
    \item \textbf{写完成标记}：非缓存写请求发出后设置 \verb|uncached_write_pending|
    \item \textbf{阻塞后续请求}：有未完成的非缓存写时，阻塞所有非缓存读写请求
    \item \textbf{bvalid 清除}：收到 AXI 写响应后清除标记，允许后续请求
    \item \textbf{Cache 行不阻塞}：Cache 行访问（type=3'b100）不受此限制
\end{itemize}

\subsubsection{Load/Store 指令的字节处理}

\begin{lstlisting}[language=Verilog, caption=Load/Store 字节对齐处理]
// Store 字节使能生成（EX 级）
wire [1:0] store_offset = data_vaddr[1:0];

always @(*) begin
    case ({is_store_byte, is_store_half, is_store_word})
        3'b100: begin  // SB
            case (store_offset)
                2'b00: data_sram_wstrb = 4'b0001;
                2'b01: data_sram_wstrb = 4'b0010;
                2'b10: data_sram_wstrb = 4'b0100;
                2'b11: data_sram_wstrb = 4'b1000;
            endcase
            data_sram_wdata = {4{ex_rkd_value[7:0]}};
        end
        3'b010: begin  // SH
            case (store_offset[1])
                1'b0: data_sram_wstrb = 4'b0011;
                1'b1: data_sram_wstrb = 4'b1100;
            endcase
            data_sram_wdata = {2{ex_rkd_value[15:0]}};
        end
        3'b001: begin  // SW
            data_sram_wstrb = 4'b1111;
            data_sram_wdata = ex_rkd_value;
        end
        default: begin
            data_sram_wstrb = 4'b0000;
            data_sram_wdata = 32'h0;
        end
    endcase
end

// Load 数据处理（MEM 级）
wire [1:0] load_offset = mem_vaddr[1:0];
wire [31:0] load_data = data_sram_rdata;

always @(*) begin
    case ({is_load_byte, is_load_half, is_load_word, 
           is_load_byte_u, is_load_half_u})
        5'b10000: begin  // LB
            case (load_offset)
                2'b00: mem_final_result = {{24{load_data[ 7]}}, load_data[ 7:0]};
                2'b01: mem_final_result = {{24{load_data[15]}}, load_data[15:8]};
                2'b10: mem_final_result = {{24{load_data[23]}}, load_data[23:16]};
                2'b11: mem_final_result = {{24{load_data[31]}}, load_data[31:24]};
            endcase
        end
        5'b01000: begin  // LH
            case (load_offset[1])
                1'b0: mem_final_result = {{16{load_data[15]}}, load_data[15:0]};
                1'b1: mem_final_result = {{16{load_data[31]}}, load_data[31:16]};
            endcase
        end
        5'b00100: begin  // LW
            mem_final_result = load_data;
        end
        5'b00010: begin  // LBU
            case (load_offset)
                2'b00: mem_final_result = {24'h0, load_data[ 7:0]};
                2'b01: mem_final_result = {24'h0, load_data[15:8]};
                2'b10: mem_final_result = {24'h0, load_data[23:16]};
                2'b11: mem_final_result = {24'h0, load_data[31:24]};
            endcase
        end
        5'b00001: begin  // LHU
            case (load_offset[1])
                1'b0: mem_final_result = {16'h0, load_data[15:0]};
                1'b1: mem_final_result = {16'h0, load_data[31:16]};
            endcase
        end
        default: mem_final_result = 32'h0;
    endcase
end
\end{lstlisting}

字节处理说明：
\begin{itemize}
    \item \textbf{Store 数据复制}：SB 将字节复制 4 份，SH 复制 2 份，便于 Cache Bank 写入
    \item \textbf{Store 字节使能}：根据地址低位生成对应的字节写使能
    \item \textbf{Load 数据选择}：根据地址低位从 32 位数据中选择对应字节或半字
    \item \textbf{符号扩展}：LB/LH 进行符号扩展，LBU/LHU 进行零扩展
\end{itemize}
